\chapter{Finalidad y Objetivos}
\label{ch:finalidad_objetivos}

En este capítulo se define el propósito estratégico del plan de seguridad digital y los objetivos medibles formulados para la UNAP.

\section{Finalidad}

La finalidad del presente Plan Institucional de Ciberseguridad es establecer, normar y articular un conjunto coordinado de medidas técnicas, organizacionales, normativas y procedimentales orientadas a prevenir, identificar, proteger, detectar, responder y recuperarse eficazmente frente a incidentes, amenazas y ataques informáticos que atenten contra la confidencialidad, integridad y disponibilidad del ecosistema tecnológico de la Universidad Nacional del Altiplano de Puno.

Para alcanzar dicha finalidad, el plan formaliza una arquitectura de defensa en profundidad mediante la sinergia de marcos y metodologías de referencia internacional:

\begin{itemize}[itemsep=3pt,topsep=2pt,parsep=0pt]
    \item \textbf{Gobernanza Estratégica (NIST CSF 2.0):} Institucionalizar la ciberseguridad a nivel transversal en la universidad, estructurando la toma de decisiones, la asignación de recursos y la gestión integral de riesgos tecnológicos dentro de las seis funciones directivas: Gobernar, Identificar, Proteger, Detectar, Responder y Recuperar.

    \item \textbf{Modelado Preventivo de Software (STRIDE):} Analizar y mitigar sistemáticamente los vectores de Suplantación (\textit{Spoofing}), Manipulación (\textit{Tampering}), Repudio (\textit{Repudiation}), Fuga de Información (\textit{Information Disclosure}), Denegación de Servicio (\textit{Denial of Service}) y Elevación de Privilegios (\textit{Elevation of Privilege}) en las aplicaciones críticas desarrolladas y operadas por la OTI (con especial énfasis en el flujo \path{matricula.unap.edu.pe}).

    \item \textbf{Vigilancia Proactiva (Threat Intelligence):} Monitorizar de forma continua indicadores de compromiso (IoCs), registros de suplantación de dominio (\textit{typosquatting}), tácticas dirigidas contra el sector público peruano y vulnerabilidades de día cero para anticiparse a los vectores de intrusión.

    \item \textbf{Caracterización Técnica del Adversario (MITRE ATT\&CK):} Mapear el comportamiento real y las tácticas, técnicas y procedimientos (TTPs) de los atacantes, habilitando capacidades analíticas en registros de auditoría y telemetría de red para transformar los incidentes en reglas de detección temprana en herramientas SIEM.

    \item \textbf{Ruptura del Vector de Intrusión (Cyber Kill Chain):} Diseñar controles específicos y puntos de quiebre defensivo en cada uno de los siete eslabones lineales del ataque, impidiendo que una amenaza en fase de reconocimiento o entrega (como el phishing facilitado por la ausencia de DMARC) culmine en la explotación o el secuestro de información sensible.

    \item \textbf{Controles Operativos Estrictos (Zero Trust Architecture):} Reemplazar la confianza implícita por una verificación continua y granular sobre identidades (MFA obligatorio), dispositivos de acceso, segmentación lógica de redes (VLANs), aplicaciones web blindadas y cifrado mandatorio de datos académicos y financieros, tanto en reposo como en tránsito.
\end{itemize}

\clearpage
\enlargethispage{2\baselineskip}
\section{Objetivos del Plan}

\subsection{Objetivo General}

Fortalecer de forma integral la postura de seguridad de la información y la resiliencia operativa de la Universidad Nacional del Altiplano de Puno (UNAP), mediante la identificación rigurosa de activos críticos, el modelado sistemático de amenazas, la implementación de controles defensivos bajo el estándar de Confianza Cero (Zero Trust) y el establecimiento de capacidades institucionales estandarizadas de detección, respuesta y recuperación continua articuladas bajo el marco NIST CSF 2.0.

\subsection{Objetivos Específicos}

\begin{enumerate}[itemsep=2.5pt,topsep=2pt,parsep=0pt]
    \item \textbf{Identificar y clasificar los activos de información:} Catalogar exhaustivamente el 100\% de los recursos tecnológicos, servidores, bases de datos, redes y servicios web asociados al dominio \path{unap.edu.pe}, estratificando su nivel de criticidad (Crítico, Alto, Medio, Bajo) conforme al impacto sobre la continuidad del servicio educativo y la confidencialidad de los datos.

    \item \textbf{Determinar amenazas y vulnerabilidades arquitectónicas:} Ejecutar el modelado preventivo de amenazas mediante la metodología STRIDE sobre los sistemas críticos de la universidad (priorizando el flujo de matrícula y campus virtual), identificando debilidades de configuración, obsolescencia de software (como \path{nginx/1.16.1}) y ausencia de directivas de protección en el transporte (\texttt{Strict-\allowbreak Transport-\allowbreak Security}).

    \item \textbf{Analizar y correlacionar las técnicas empleadas por atacantes:} Caracterizar el comportamiento potencial de ciberadversarios en el contexto universitario mediante la matriz MITRE ATT\&CK Enterprise, traduciendo las técnicas y procedimientos observados en reglas formales de visibilidad, auditoría y monitoreo de eventos.

    \item \textbf{Detectar amenazas reales mediante inteligencia proactiva:} Institucionalizar un flujo operativo de Threat Intelligence para monitorizar activamente campañas de Spear-Phishing, registros fraudulentos de dominios (\textit{typosquatting}), abuso de la reputación institucional derivado de la ausencia de DMARC y explotación de vulnerabilidades conocidas (CVEs).

    \item \textbf{Implementar controles preventivos y perimetrales:} Desplegar mecanismos de mitigación perimetral (Web Application Firewall - WAF, inspección de cabeceras, políticas DNS estrictas SPF: \texttt{-all} y DMARC: \texttt{p=reject}) y capacidades de protección en estaciones de trabajo y servidores mediante soluciones EDR centralizadas.

    \item \textbf{Fortalecer el gobierno de identidades y el control de accesos:} Implementar la directriz de Autenticación Multifactor (MFA) obligatoria para personal administrativo, directivo y docente, combinada con políticas de privilegios mínimos (\textit{Principle of Least Privilege}) y asignación estricta de roles (RBAC) en las plataformas académicas.

    \item \textbf{Reducir los riesgos de exposición y fuga de datos personales:} Blindar los canales de transmisión y repositorios que albergan calificaciones, historiales académicos y datos de identificación personal (DNI, registros biométricos), garantizando el cumplimiento irrestricto de la Ley N° 29733 (Ley de Protección de Datos Personales).

    \item \textbf{Mejorar la capacidad de respuesta y resiliencia ante incidentes:} Formalizar el Equipo de Respuesta a Incidentes de Seguridad Digital de la UNAP (CSIRT-UNAP), definiendo procedimientos operativos estandarizados (\textit{Playbooks}) para minimizar los tiempos medios de detección (MTTD) y contención (MTTR), asegurando la restauración íntegra de servicios mediante copias de seguridad inmutables (estrategia 3-2-1).
\end{enumerate}
